% !TeX root = ../main.tex

\chapter{Networking}
\label{chap:Networking}
\section{Overview}
The default configuration for these devices is to use static IP addresses on the 10.69.26.0/23 subnet.\\\\
The network topology is generally as follows: 
\begin{figure}[htbp]
	\centering
	\begin{tikzpicture}[
		node/.style={draw, rectangle, minimum width=3.5cm, minimum height=1.2cm, align=center},
		switch/.style={draw, rectangle, minimum width=4cm, minimum height=1.5cm, fill=gray!10, align=center},
		line/.style={-latex, thick}
		]
		
		% ======================
		% Network Switch
		% ======================
		\node[switch] (sw) at (0,0) {Network Switch};
		
		% ======================
		% LEFT SIDE
		% ======================
		\node[node] (workstation) at (-6,1.5) {Workstation\\IP Address: xxx};
		
		\node[node] (ate) at (-6,-0.5) {ATE\\IP Address: xxx};
		
		\node[node] (fofb1) at (-6,4.8) {PSC FOFB \#1\\IP Address: xxx};
		
		% ======================
		% RIGHT SIDE
		% ======================
		\node[node] (pscnet) at (6,1.5) {PSC Network\\IP Address: xxx};
		
		\node[node] (fofb2) at (6,0) {PSC FOFB \#2\\IP Address: xxx};
		
		% ======================
		% CONNECTIONS
		% ======================
		\draw[line] (workstation) -- (sw);
		\draw[line] (ate) -- (sw);
		\draw[line] (pscnet) -- (sw);
		\draw[line] (fofb2) -- (sw);
		
		% ======================
		% DIRECT CONNECTION
		% ======================
		\draw[line, dashed, red, thick] (fofb1) -- (workstation);
		
		\node[red, font=\small, align=center]
		at ($(fofb1)!0.5!(workstation) + (2.6,0.4)$)
		{Direct connection\\FOFB \#1 $\leftrightarrow$ Workstation};
		
	\end{tikzpicture}
	\caption{Network Topology with Clean FOFB Direct Link}
\end{figure}